\chapter{Public Goods and Common Resources}
\label{chap:public-goods}

\section{Learning objectives}

After completing this chapter, you should be able to classify goods by rivalry and excludability, explain free riding and common-resource overuse, and evaluate institutional responses.

\section{Classifying goods}

A good is \keyterm{rival} if one person's consumption reduces what remains for others, and \keyterm{excludable} if non-payers can be prevented from consuming it. Private goods are rival and excludable. \keyterm{Public goods} are non-rival and non-excludable, such as national defence. \keyterm{Common resources} are rival but non-excludable, such as an open fishery. Club goods are non-rival up to congestion but excludable.

\section{Free riding and the commons}

Non-excludability encourages \keyterm{free riding}: people may benefit without contributing, leading private provision of public goods to be too low. Common resources face the opposite problem: individuals have an incentive to use the resource before others do, causing overuse, the \keyterm{tragedy of the commons}.

Responses include public provision, taxation, user fees, property rights, quotas, community governance, and rules that monitor and sanction use. The appropriate institution depends on information, enforcement capacity, values, and local conditions.

\section*{Chapter Review}

\subsection*{Key concepts}
\begin{itemize}
  \item \textbf{Rivalry}: one person's use reduces availability to others.
  \item \textbf{Excludability}: non-payers can be prevented from consuming.
  \item \textbf{Free riding}: benefiting without contributing.
  \item \textbf{Tragedy of the commons}: overuse of a non-excludable rival resource.
\end{itemize}

\section*{Exercises}

\subsection*{Conceptual questions}
\begin{enumerate}[label=\arabic*.]
  \item Classify a public park, streaming subscription, fishery, and sandwich by rivalry and excludability.
  \item Why does free riding lead to underprovision of public goods?
  \item Why is a common resource vulnerable to overuse?
\end{enumerate}

\subsection*{Graphical exercises}
\begin{enumerate}[resume,label=\arabic*.]
  \item Draw a simple diagram showing private and social benefit for a public good.
\end{enumerate}

\subsection*{Analytical exercises}
\begin{enumerate}[resume,label=\arabic*.]
  \item Explain why assigning transferable fishing quotas can reduce overfishing.
  \item Compare public provision and community governance of a local park.
\end{enumerate}

\subsection*{Exam-style multiple-choice questions}
\begin{enumerate}[resume,label=\arabic*.]
  \item A public good is: (a) rival and excludable; (b) non-rival and non-excludable; (c) rival and non-excludable; or (d) always supplied by government.
  \item A common resource is typically: (a) non-rival and excludable; (b) rival and non-excludable; (c) non-rival and non-excludable; or (d) privately owned by definition.
\end{enumerate}

\subsection*{PPE discussion question}

\noindent Is private property always the best solution to common-resource problems? Discuss ecological, distributive, and institutional considerations.
